% Lean module: MonadFormal.Refinement
\subsubsection{Semantic refinement versus obtainable refinement}

Let $q:X\to Y$ be retained information and $r:X\to Z$ a newly
required distinction. Define $q^+(x)=(q(x),r(x))$.
\leandef{MonadFormal.refine}

\begin{theorem}[Least common refinement]\label{thm:least-refinement}
The map $q^+$ preserves both $q$ and $r$, and
\[
 q^+(x)=q^+(y)\iff q(x)=q(y)\ \text{and}\ r(x)=r(y).
\]
If both $q$ and $r$ factor through $s:X\to W$, then so does $q^+$.
Thus $q^+$ is least in the order determined by factorization.
\end{theorem}
\begin{proof}
The coordinate projections recover $q$ and $r$. Equality of ordered
pairs is equality of both coordinates. If $q=g\circ s$ and
$r=h\circ s$, then $q^+=(g,h)\circ s$.
\leanproof{MonadFormal.refine_fibers,MonadFormal.refine_preserves,MonadFormal.refine_least}
\end{proof}

\begin{theorem}[Obtainability from declared access]\label{thm:obtainable-refinement}
Let $A:X\to W$ describe the information actually accessible to the
system. The refined value $q^+$ factors through $A$ if and only if
both $q$ and $r$ factor through $A$. Moreover, deterministic
postprocessing of $A$ cannot separate any pair that $A$ identifies.
\end{theorem}
\begin{proof}
The forward implication composes the two projections with a decoder
for $q^+$. The reverse implication is the least-refinement property.
For the last statement, $A(x)=A(y)$ implies $h(A(x))=h(A(y))$ for
every function $h$.
\leanproof{MonadFormal.factors_trans,MonadFormal.refine_obtainable_iff,MonadFormal.processing_preserves_fibers}
\end{proof}
The definition of $q^+$ supplies no access to $r(x)$. In particular,
the analyst's access to $x$ is not an admissible operation of the system.

\begin{example}[New carrier states need not split old states]
\label{ex:carrier-extension}
Let $X=\{*\}$, $X'=\{*\}\times\{0,1\}$ and $e(*)=(*,0)$.
The second-coordinate probe is constant after pullback to $X$, so
it factors through the constant old representation. Nevertheless,
it separates $(*,0)$ and $(*,1)$ in $X'$, and $(*,1)$ is outside
the image of $e$.
\end{example}
\begin{proof}
The pulled-back probe has value zero. The two displayed target values
are zero and one; no image point has second coordinate one.
\leanproof{MonadFormal.extra_bit_pullback_factors,MonadFormal.extra_bit_separates_new_states,MonadFormal.extra_state_outside_image}
\end{proof}
Pullback sufficiency over old states therefore does not describe all
distinctions of an enlarged carrier.
